\subsection{KT-TWOFIELD-661: simultaneous genuine Eisenstein congruences}
\label{sec:kt-twofield-661}

Put $\chi_D=\chi_{-D}$ for $D=3,7$, and use the standard
normalization
\[
 \sum_{a=1}^{D}\frac{\chi_D(a)t e^{at}}{e^{Dt}-1}
       =\sum_{j\geq0}B_{j,\chi_D}\frac{t^j}{j!}.
                                                               \tag{TF.1}
\]
The following example gives a concrete obstruction to proving
an all-weight unit theorem merely by switching between the
two classical auxiliary levels and varying split Frobenius
primes. It is not an obstruction to the rational $K$-group
vanishing, which the separate reflected test actually proves
at this prime.

\begin{proposition}[KT-TWOFIELD-661]
Let $e=268$, $w=2e+1=537$, and $\ell=661$. Then
\[
 \ell\mid\operatorname{num}\frac{B_{w,\chi_3}}w,
 \qquad
 \ell\mid\operatorname{num}\frac{B_{w,\chi_7}}w,
 \qquad
 B_{\ell-1-e}=B_{392}\equiv53\pmod\ell.
                                                               \tag{TF.2}
\]
For each $D=3,7$ there are a normalized primitive newform
$f_D\in S_{537}(\Gamma_0(D),\chi_D)$ and a coefficient
prime $\lambda_D\mid661$ for which
\[
 a_q(f_D)\equiv1+\chi_D(q)q^{536}\pmod{\lambda_D}
                  \qquad(q\text{ prime},\ q\nmid D).
                                                               \tag{TF.3}
\]
Consequently every split-prime classical characteristic
polynomial comparison at $1+q^{536}$ is divisible by $661$
at both auxiliary levels. In the full level-seven model,
the twist-inclusive determinant
$\det((1+2^{536})^3I-T_2^3)$ is also nonunit at $661$.
Nevertheless the carried K17-REFLECTOR gives
\[
                         K_{536}(\mathbf Z)\{661\}=0.
                                                               \tag{TF.4}
\]
\end{proposition}

\begin{proof}
The prime $661$ is proved prime by trial division through
$25$. It satisfies $661>538$ and is coprime to $42\cdot537$.
The complete standard-library certificate
\texttt{bernoulli\_661\_certificate.py} uses two independent
finite identities to verify the generalized residues.

First, the defining recurrence
\[
 B_0=1,\qquad
 B_m=-\frac1{m+1}\sum_{j=0}^{m-1}\binom{m+1}{j}B_j
                                                               \tag{TF.5}
\]
is computed modulo $661$ through $m=537$. Every divisor
$m+1$ used is smaller than $661$, so it is a unit.
Bernoulli polynomials and (TF.1) give the finite identity
\[
 B_{w,\chi_D}
   =\sum_{a=1}^{D}\chi_D(a)
     \sum_{j=0}^{w}\binom wj B_j a^{w-j}D^{j-1}.
                                                               \tag{TF.6}
\]
Its residues are zero for both $D$; division by $w$ is
permitted. The same recurrence gives $B_{392}\equiv53$.

As an independent check of the generalized values, set
$S_D=\sum_{a=1}^{D\ell}\chi_D(a)a^w$.
Multiplication of (TF.1) by $(e^{D\ell t}-1)/t$ gives
\[
 S_D=\sum_{j=0}^{w}\binom wj B_{j,\chi_D}
             \frac{(D\ell)^{w+1-j}}{w+1-j}
       \equiv D\ell B_{w,\chi_D}\pmod{\ell^2}.
                                                               \tag{TF.7}
\]
All denominators are units, and the $B_{j,\chi_D}$ are
$\ell$-integral by (TF.5)--(TF.6). Direct modular exponentiation
gives $S_3\equiv S_7\equiv0\pmod{661^2}$, again proving
the first two congruences. A separate FLINT exact-rational
calculation, stored with complete numerators, gives
valuation one in both cases, with respective denominators
$9$ and $7$ for $B_{537,\chi_D}/537$. Exact valuation is
not needed for the conclusion.

Now apply Fretwell--Roberts, \emph{Newform Eisenstein
congruences of local origin}, Ramanujan J.~64 (2024),
505--527, Theorem~3.6, pp.~517--518, with
\[
 N=D,\quad M=1,\quad k=537,\quad
 \psi=1,\quad\phi=\chi_D,\quad\lambda'=(661).
\]
The conductor is squarefree, $k>2$, and $661\nmid6D$.
The required positive valuation is precisely that of
$L(1-k,\chi_D)=-B_{k,\chi_D}/k$, established above.
The theorem supplies (TF.3); its paragraph after the proof
states that $M=1$ and primitive character make the form new.
The theorem's entire displayed proof was read in the original
published source. Its modular reduction and lifting inputs
are cited, not independently re-proved here.

For a split prime, $\chi_D(q)=1$, so (TF.3) makes
$1+q^{536}-a_q(f_D)$ nonunit. The characteristic polynomial
of the full classical space contains this eigenvalue,
and its evaluation is therefore divisible by $661$.
For level seven the quadratic-character space is contained
in the full $\Gamma_1(7)$ space. Since $\chi_7(2)=1$,
the eigenvalue $a_2(f_7)$ makes the twist-inclusive cube
determinant nonunit as asserted.

Finally $2\leq e=268\leq658=\ell-3$ and the nonzero
residue $B_{392}\equiv53$ meets exactly the hypotheses
of K17-REFLECTOR, (RB.1), K17.tex lines~4533--4619.
That carried theorem kills the cyclotomic even component
$A_{661}^{[-268]}=A_{661}^{[392]}$, hence (TF.4).
\end{proof}

\paragraph{The precise all-weight gate exposed by the example.}
For even $e\geq2$ define
\[
 \mathcal E_D(e)=\left\{\ell>2e+2:\ \ell\nmid6D,\quad
       \ell\mid\operatorname{num}
                 \frac{B_{2e+1,\chi_D}}{2e+1}\right\}.
                                                               \tag{TF.8}
\]
Theorems~3.5--3.6 of Fretwell--Roberts, in their $M=1$
case, identify this with the prime support of the indicated
classical Eisenstein eigensystem congruences. These are
simultaneous congruences at all good primes, stronger than
an accidental common divisor of two determinant evaluations.
Equation (TF.2) disproves the hoped-for assertion
$\mathcal E_3(e)\cap\mathcal E_7(e)=\varnothing$.

A sufficient arithmetic statement for eliminating this
particular common support would instead be
\[
 \ell\in\mathcal E_3(e)\cap\mathcal E_7(e)
          \quad\Longrightarrow\quad
                   \ell\nmid B_{\ell-1-e}.
                                                               \tag{TF.9}
\]
It holds in the displayed example. No theorem establishing
(TF.9) uniformly has been obtained in this run. Moreover,
(TF.9) alone would not complete the full Selmer argument:
the independent rank-three unitary comparison can also be
nonunit, even when a classical branch is a unit. The
complete all-weight problem requires its actual joint
exceptional support to avoid the rational component
$A_\ell^{[-e]}$. A generalized Bernoulli numerator attached
to $\chi_D$ cannot simply be identified with this component.
The reflected ordinary Bernoulli number in (TF.9) records
the exact distinct character required by the supplied
class-field argument.

One can also see the distinction directly in the residual
representations. The eigensystem in (TF.3) has semisimplified
characters $1$ and $\chi_D\omega^{2e}$; after the weight-zero
twist these become $\omega^{-e}$ and $\chi_D\omega^e$.
An extension between them has ratio character
$\chi_D\omega^{2e}$ or $\chi_D\omega^{-2e}$, whereas
the rational negative Selmer problem uses $\omega^{-e}$.
Restriction to $G_{\mathbf Q(\sqrt{-D})}$ removes
$\chi_D$ but does not remove the doubling of the exponent.
For odd $\ell$, restriction of a rational cohomology class
to this quadratic field is injective by restriction--corestriction,
whose composite is multiplication by two. Thus field
variation supplies a new automorphic comparison, not a
formal annihilation of the original rational class.

